%!TeX root=../main.tex

\chapter{Desenvolvimento}
\label{cap:desenvolvimento}

\section{Pixelização}
\label{sec:pixelizacao}

Para alcançar a estética desejada, a pixelização do mundo não foi
implementada por meio de um efeito de pós-processamento aplicado sobre a
imagem final, mas sim por redução direta de escala (\textit{downscale}) no \textit{viewport}.
Em vez de renderizar a cena em resolução superior e executar um cálculo
adicional de pós-processamento apenas para reduzir a qualidade visual,
renderizar diretamente em uma resolução interna menor evita trabalho
desnecessário de rasterização e economiza processamento da GPU.

O projeto adota uma resolução base de referência de $1600 \times 1200$ com o modo
de `esticar' configurado em \textbf{\textit{keep\_height}} na Godot, mantendo a altura de $1200$ como
referência vertical e ajustando a largura proporcionalmente ao \textit{aspect
ratio} da tela (como monitores 21:9). O sistema de renderização foi estruturado a
partir do nó \textbf{MainGame} com a separação do fluxo visual em diferentes \textit{viewports}.
Essa divisão permite aplicar o efeito de pixelização exclusivamente ao espaço
do mundo, preservando a nitidez da interface e das instruções na cena.

O primeiro \textit{viewport} corresponde ao \textbf{GameViewport} --- que é um nó \textit{SubViewportContainer} e seu respectivo
\textit{SubViewport} que abriga o mundo do jogo --- e executa a redução de escala
(\textit{downscale}) por meio do parâmetro `taxa de pixelização', que controla o
\textbf{\textit{stretch\_shrink}} do \textit{viewport}. Com a taxa padrão igual a \textbf{4}, o \textbf{\textit{stretch\_shrink}}
assume o valor 4 e a renderização interna do mundo é reduzida
proporcionalmente (de $1600 \times 1200$ para $400 \times 300$ na resolução base), sendo
reescalonada para a tela com aspecto pixelado.

O segundo \textit{subviewport} é dedicado a elementos auxiliares de HUD que
aparecem na cena --- como indicadores flutuantes sobre itens interativos
(por exemplo, quando o jogador se aproxima do item \textbf{Tampa} para coletá-lo,
é exibida a tecla \textit{E} se ele estiver jogando no teclado.) ---, permanecendo fora do
downscale do \textbf{GameViewport}. Já telas como o menu principal são compostas em camadas
de interface sobrepostas posteriormente, sem depender desse \textit{subviewport} de cena
nem do uso de \textit{spatial shaders} voltados ao espaço 3D, o que garante leitura nítida
dos menus.

Essa arquitetura garante que a estética pixelizada do mundo não atrapalhe
a legibilidade das interfaces e permite que sejam ajustáveis. Porém, essa abordagem
acaba gerando circunstâncias que influenciam o cálculo de entradas e a rotação da
câmera, detalhados na próxima seção.